% GENERATED FILE. Edit canonical lessons, then run npm run generate:books.
% canonical-source-hash: fnv1a64:5a6fe72384c958eb
% canonical-lessons: GE-C346-zweck, GE-C346-absicht, GE-C346-damit, GE-C346-absichtlich, GE-C346-anlass

\chapter{Wozu?}
\label{ch:ge-wozu}
\hlchaptermodality{346}

\begin{chapteropening}
\textbf{By the end of this chapter:} \emph{I can say what something is for, what someone intends, and name the aim of an action with damit.}

Complete the last lesson of chapter 346: I can say what something is for, what someone intends, and name the aim of an action with damit.
\end{chapteropening}

\section[der Zweck]{der Zweck — the purpose}

\label{lesson:GE-C346-zweck}



\begin{quote}
Before the new one: say the German for to back up, then the German for to prove.
\end{quote}

\subsection*{You'll want to know: der Zweck}
\textbf{der Zweck} — \textquotedblleft{}the purpose\textquotedblright{}.

\begin{quote}
Der Zweck der Reise war ein Besuch bei meiner Tante.
\end{quote}

\begin{quote}
\emph{The purpose of the trip was a visit to my aunt.}
\end{quote}

Plural: \emph{die Zwecke}. \emph{Das hat keinen Zweck} — there is no point.

\begin{cousinweb}
\emph{Zweck} is inherited from Old High German \emph{zwec} (EtymDB), a word for a wooden peg or nail. It is usually explained as the peg in the middle of a target: the thing you aim at.
\end{cousinweb}

\begin{grammarlens}[title={What you've built}]
One more word for telling, explaining and describing at B1.
\end{grammarlens}

\subsection*{Your turn}
\begin{itemize}
\raggedright
  \item \emph{Say it:} \emph{der Zweck}
  \item \emph{Say it:} \emph{der Zweck}, once more
  \item \emph{Say it:} a sentence of your own with \emph{der Zweck}
  \item \emph{Recall:} say the German for to back up, then the German for to prove, then say \emph{der Zweck} again
\end{itemize}

\begin{tcolorbox}[breakable,colback=teal!4,colframe=teal!35!black,title={Before you move on}]
What does der Zweck mean? (\textquotedblleft{}The purpose\textquotedblright{}.) Say it once more.
\end{tcolorbox}
\section[die Absicht]{die Absicht — the intention}

\label{lesson:GE-C346-absicht}



\begin{quote}
Before the new one: say the German for to prove, then the German for the purpose.
\end{quote}

\subsection*{You'll want to know: die Absicht}
\textbf{die Absicht} — \textquotedblleft{}the intention\textquotedblright{}.

\begin{quote}
Es war nicht meine Absicht, so laut zu sein.
\end{quote}

\begin{quote}
\emph{It was not my intention to be so loud.}
\end{quote}

Plural: \emph{die Absichten}. \emph{Mit Absicht} — on purpose.

\begin{cousinweb}
\emph{Absicht} is built from \emph{ab} and \emph{Sicht} \textquotedblleft{}sight\textquotedblright{}, from the family of \emph{sehen} \textquotedblleft{}to see\textquotedblright{}, Old High German \emph{sehan} (EtymDB): the thing you have your eye on.
\end{cousinweb}

\begin{grammarlens}[title={What you've built}]
One more word for telling, explaining and describing at B1.
\end{grammarlens}

\subsection*{Your turn}
\begin{itemize}
\raggedright
  \item \emph{Say it:} \emph{die Absicht}
  \item \emph{Say it:} \emph{die Absicht}, once more
  \item \emph{Say it:} a sentence of your own with \emph{die Absicht}
  \item \emph{Recall:} say the German for to prove, then the German for the purpose, then say \emph{die Absicht} again
\end{itemize}

\begin{tcolorbox}[breakable,colback=teal!4,colframe=teal!35!black,title={Before you move on}]
What does die Absicht mean? (\textquotedblleft{}The intention\textquotedblright{}.) Say it once more.
\end{tcolorbox}
\section[damit]{damit — in order that, so that; with it}

\label{lesson:GE-C346-damit}



\begin{quote}
Before the new one: say the German for the purpose, then the German for the intention.
\end{quote}

\subsection*{You'll want to know: damit}
\textbf{damit} — \textquotedblleft{}in order that, so that; with it\textquotedblright{}.

\begin{quote}
Ich spreche langsam, damit alle mich verstehen.
\end{quote}

\begin{quote}
\emph{I speak slowly so that everyone understands me.}
\end{quote}

As a conjunction \emph{damit} sends the verb to the end and names an aim. As an adverb it means with it: \emph{Was machst du damit?}

\begin{cousinweb}
\emph{damit} joins \emph{da} \textquotedblleft{}there\textquotedblright{} and \emph{mit} \textquotedblleft{}with\textquotedblright{}, from Old High German \emph{miti} (EtymDB). It first meant \textquotedblleft{}with that\textquotedblright{}; the sense \textquotedblleft{}so that\textquotedblright{} grew out of it.
\end{cousinweb}

\begin{grammarlens}[title={What you've built}]
One more word for telling, explaining and describing at B1.
\end{grammarlens}

\subsection*{Your turn}
\begin{itemize}
\raggedright
  \item \emph{Say it:} \emph{damit}
  \item \emph{Say it:} \emph{damit}, once more
  \item \emph{Say it:} a sentence of your own with \emph{damit}
  \item \emph{Recall:} say the German for the purpose, then the German for the intention, then say \emph{damit} again
\end{itemize}

\begin{tcolorbox}[breakable,colback=teal!4,colframe=teal!35!black,title={Before you move on}]
What does damit mean? (\textquotedblleft{}In order that, so that; with it\textquotedblright{}.) Say it once more.
\end{tcolorbox}
\section[absichtlich]{absichtlich — on purpose, deliberate}

\label{lesson:GE-C346-absichtlich}



\begin{quote}
Before the new one: say the German for the intention, then the German for in order that.
\end{quote}

\subsection*{You'll want to know: absichtlich}
\textbf{absichtlich} — \textquotedblleft{}on purpose, deliberate\textquotedblright{}.

\begin{quote}
Er hat das Glas nicht absichtlich fallen lassen.
\end{quote}

\begin{quote}
\emph{He did not drop the glass on purpose.}
\end{quote}

The opposite is \emph{aus Versehen} — by mistake.

\begin{cousinweb}
\emph{absichtlich} is built with \emph{-lich} on \emph{Absicht} \textquotedblleft{}intention\textquotedblright{}, from \emph{ab} and \emph{Sicht} \textquotedblleft{}sight\textquotedblright{}, of the family of \emph{sehen} \textquotedblleft{}to see\textquotedblright{} (EtymDB).
\end{cousinweb}

\begin{grammarlens}[title={What you've built}]
One more word for telling, explaining and describing at B1.
\end{grammarlens}

\subsection*{Your turn}
\begin{itemize}
\raggedright
  \item \emph{Say it:} \emph{absichtlich}
  \item \emph{Say it:} \emph{absichtlich}, once more
  \item \emph{Say it:} a sentence of your own with \emph{absichtlich}
  \item \emph{Recall:} say the German for the intention, then the German for in order that, then say \emph{absichtlich} again
\end{itemize}

\begin{tcolorbox}[breakable,colback=teal!4,colframe=teal!35!black,title={Before you move on}]
What does absichtlich mean? (\textquotedblleft{}On purpose, deliberate\textquotedblright{}.) Say it once more.
\end{tcolorbox}
\section[der Anlass]{der Anlass — the occasion; the cause}

\label{lesson:GE-C346-anlass}



\begin{quote}
Before the new one: say the German for in order that, then the German for on purpose.
\end{quote}

\subsection*{You'll want to know: der Anlass}
\textbf{der Anlass} — \textquotedblleft{}the occasion; the cause\textquotedblright{}.

\begin{quote}
Zu diesem Anlass trägt er einen Anzug.
\end{quote}

\begin{quote}
\emph{For this occasion he is wearing a suit.}
\end{quote}

Plural: \emph{die Anlässe}. \emph{Aus diesem Anlass} — on this occasion, for this reason. In Switzerland \emph{der Anlass} is also an event.

\begin{cousinweb}
\emph{Anlass} is built from \emph{an} and \emph{lassen} \textquotedblleft{}to let, to let go\textquotedblright{}, inherited from Old High German \emph{lāzan} and Proto-Germanic \emph{lētaną} (EtymDB): something that lets events start.
\end{cousinweb}

\begin{grammarlens}[title={What you've built}]
One more word for telling, explaining and describing at B1.
\end{grammarlens}

\subsection*{Your turn}
\begin{itemize}
\raggedright
  \item \emph{Say it:} \emph{der Anlass}
  \item \emph{Say it:} \emph{der Anlass}, once more
  \item \emph{Say it:} a sentence of your own with \emph{der Anlass}
  \item \emph{Recall:} say the German for in order that, then the German for on purpose, then say \emph{der Anlass} again
\end{itemize}

\begin{tcolorbox}[breakable,colback=teal!4,colframe=teal!35!black,title={Before you move on}]
What does der Anlass mean? (\textquotedblleft{}The occasion; the cause\textquotedblright{}.) Say it once more.
\end{tcolorbox}
